% Submission template — ICLR-style section skeleton.
%
% NOTE ON THE DOCUMENT CLASS: ICLR's official iclr2027_conference.sty is not yet
% published and is not a CTAN package, so this template uses `article` with
% settings that approximate the target layout. Swapping in the official style
% file is a mechanical change once it is released: replace the \documentclass
% and geometry lines, nothing else.
\documentclass[10pt,twocolumn]{article}

\usepackage[letterpaper,margin=0.9in,columnsep=0.28in]{geometry}
\usepackage[T1]{fontenc}
\usepackage[utf8]{inputenc}
\usepackage{times}
\usepackage{amsmath}
\usepackage{amssymb}
\usepackage{graphicx}
\usepackage{booktabs}
\usepackage{multirow}
\usepackage{caption}
\usepackage[hidelinks]{hyperref}
\usepackage{natbib}
\usepackage{xcolor}
\usepackage{microtype}
\usepackage{nicefrac}
\usepackage{url}
\usepackage{cleveref}
\usepackage{microtype}
\usepackage{nicefrac}
\usepackage{cleveref}

\captionsetup{font=small,labelfont=bf}
\setlength{\parskip}{2pt}
\sloppy
\emergencystretch=2em

\title{Reading Is Not Reasoning, and Reasoning Is Not Robustness:\\
Decomposing What a Language Model Does with a Solver Certificate}
\author{Anonymous authors\\Paper under double-blind review}
\date{}

\begin{document}
\maketitle

\begin{abstract}
Solver-, tool- and retrieval-augmented language model systems report large
accuracy gains, and those gains are routinely read as evidence that the model
used the supplied material. Two mechanisms produce the same number: the model
reading a supplied answer, and the model reasoning over supplied state. The
standard control, a same-shape irrelevant record, separates neither. We build a
control that does. Across three sealed experiments on frozen small models,
scored by direct next-token likelihood over verified single-token candidates, we
decompose solver assistance into an answer channel and a state channel. Reading
a completed five-arm factorial as a $2\times2$ of answer evidence by proof state
recovers a state main effect of $+22.4$pp that its prespecified diagonal
contrast could not attribute. Holding line count, line types, query-entity
occurrence count, query-predicate presence and token length fixed across ten
arms and destroying validity alone gives a validity component of $+60.4$pp
against a surface component of $+1.0$pp: $98.3\%$ of the state effect is
inference over the supplied lines, not text matching. That inference is exactly
one step deep --- withholding two steps rather than one scores identically to
supplying nothing --- and it confers no robustness: given a valid-looking
certificate whose fabricated final rule establishes the negation, the model
follows it on 189 of 192 items ($d' = -4.36$). A ten-model three-class sweep,
whose \emph{undetermined} items are certified by forward-closure saturation
validated at $23{,}240/23{,}240$ against an independent corpus, finds three
viable substrates of ten and a minimum viable interface size near 3B, not
monotonic in scale. Two measurement lessons recur: an arm scoring $50.0\%$ on a
balanced panel with $d' = 0.00$, and a prespecified criterion that a degenerate
single-label responder maximised. We report sensitivity alongside accuracy
throughout, and never upgrade a system-level effect into a claim about the
model's own reasoning.
\end{abstract}

\section{Introduction}

Pairing a language model with an external solver, tool or retrieval index
reliably raises accuracy on tasks the model cannot do alone. Systems that
translate a natural-language problem into a formal representation, dispatch it
to a symbolic engine, and return the engine's verdict through the model report
large end-to-end gains on logical reasoning
\citep{pan2023logiclm,olausson2023linc,ye2023satlm}, on arithmetic and
algorithmic tasks solved by an interpreter \citep{gao2022pal}, on API-mediated
subtasks \citep{schick2023toolformer}, on formal theorem proving
\citep{yang2023leandojo}, and on knowledge-intensive generation supported by
retrieved passages \citep{lewis2020retrievalaugmented}. Those gains are
routinely read as evidence that the model \emph{used} the supplied material.

That reading is underdetermined. At least two mechanisms produce the same
accuracy number. Under \emph{answer reading}, the supplied record contains the
conclusion, or something that co-occurs with it, and the model copies out the
implied response. Under \emph{state use}, the model works with the supplied
premises and completes the remaining inference itself. The first says nothing
about what the model does with reasoning; the second says something specific.
The standard control in this literature is a same-shape but irrelevant record,
which rules out proof length, generic validity and formatting. It does not rule
out answer reading, because a \emph{relevant} certificate stays
answer-diagnostic even when the response token is withheld: a prefix that stops
one step short still names the query's subject and predicate in adjacent lines,
so a model that only checks surface co-occurrence reproduces every number a
reasoning model would.

Separating the two requires a control that holds surface form fixed and
destroys validity alone. This paper builds one. Three sealed experiments on
frozen small models score responses by direct next-token likelihood over
verified single-token candidates, so no free text is parsed and no verdict is
inferred from prose. The first reads a completed five-arm certificate factorial
as a $2\times2$ of \emph{answer evidence} $\times$ \emph{proof state}, which
separates the two channels that end-to-end accuracy conflates. The second holds
line count, line types, query-entity occurrence count, query-predicate presence
and token length fixed across ten arms and breaks only the logic, which splits
the state effect into a surface-overlap component and a validity-tracking
component, and additionally varies entity repetition, line order and truncation
depth, and supplies a valid-looking certificate whose fabricated final rule
points the wrong way. The third scores ten models on a three-class panel whose
\emph{undetermined} items are certified by forward-closure saturation, and asks
which models can serve as the interface at all.

The picture that emerges is narrower than the headline accuracy suggests.
Holding surface form fixed and breaking only validity removes almost the entire
state effect: the validity component is $+60.4$pp against a surface component
of $+1.0$pp, so $98.3\%$ of the state effect is inference over the supplied
lines rather than text matching. Entity repetition contributes nothing --- a
valid derivation about the query entity toward an unrelated predicate scores
identically ($+0.0$pp) to one about a different entity, despite naming the query
subject four times rather than zero. The inference is exactly one step deep:
withholding two steps instead of one scores identically to supplying no
certificate at all, a cliff rather than a slope. And the tracking confers no
robustness: handed a valid-looking certificate whose fabricated final rule
establishes the negation, the model follows it on 189 of 192 items. Across ten
models, three clear a per-class recall floor; below 1.7B parameters models
collapse to a single label whatever the solver supplies, and the largest model
tested collapses while a 3B model does not.

Two measurement points shape how any of this may be reported. First, accuracy
alone is not interpretable here: one arm scored $50.0\%$ on a balanced panel
with $d' = 0.00$, matching a trivial always-one-label strategy while
discriminating nothing, so sensitivity and criterion are reported alongside
accuracy throughout \citep{macmillan1991detection,stanislaw1999calculation}.
Paired contrasts use the exact McNemar test \citep{mcnemar1947note} and
bias-corrected accelerated bootstrap intervals \citep{efron1987better}. Second,
the claim boundary is strict and is maintained everywhere in this paper. Every
arm without a near-complete supplied chain returns $d' = 0$, so the models
studied do not reason unaided; high accuracy with full solver material describes
a faithful transcriber, not a reasoner. System-level effects are never upgraded
into claims about the model's own reasoning.

The contributions are: (i) an \emph{answer-evidence decomposition} that
separates solver-supplied answers from solver-supplied reasoning state, reusable
for any tool-augmented system; (ii) a \emph{surface-matched invalidity control}
that isolates logical validity from every co-occurring surface cue, with
per-item certification that the displayed lines do not settle the query; (iii) a
\emph{forward-closure certifier} for constructing formally underdetermined
items, validated at $23{,}240/23{,}240$ against an independently labelled
corpus \citep{tafjord2020proofwriter}; (iv) evidence that validity tracking is
real, shallow, and offers no protection against corrupted solver output, a
reliability result for the whole class of solver-in-the-loop architectures; and
(v) a substrate-size floor from a ten-model sweep spanning open instruction-tuned
families \citep{team2024qwen,grattafiori2024llama,team2025gemma,microsoft2025phi,team2024olmo,allal2025smollm}.

\section{Related Work}

\subsection{Faithfulness of displayed reasoning}

Chain-of-thought prompting made intermediate reasoning steps a standard part of
the model's output \citep{wei2022chainofthought}, and a body of work then asked
whether those steps are what the answer actually rests on.
\citet{turpin2023language} show that explanations systematically fail to mention
biasing features that demonstrably drive the prediction, so a plausible trace
can misrepresent the true cause. \citet{lanham2023measuring} intervene directly
on the trace --- truncating it so the model must answer early, and injecting
mistakes into it --- and read the resulting change in the answer as a measure of
how strongly the answer is conditioned on the displayed reasoning; they find
wide variation across tasks and models, and that mistakes are often ignored.
The same manipulations have since been applied to naturally worded prompts
without artificial bias \citep{arcuschin2025chainofthought} and to models trained
to produce long reasoning traces, where hints that determine the answer are
verbalised only a minority of the time \citep{chen2025reasoning}. A separate
response to the problem is architectural: \citet{lyu2023faithful} make the trace
a formal program executed by a deterministic solver, so the displayed chain is
faithful to the answer by construction.

The limitation this paper addresses is one of object rather than of method.
This literature intervenes on \emph{model-generated} traces, where the question
is whether the model's own account is honest. It does not address
\emph{externally supplied, formally verified} state, where the material is
correct by construction and the question inverts: does the model track the
validity of a derivation it did not produce? This paper adapts the truncation
and mistake-injection manipulations to solver certificates, and adds a control
this literature does not use --- an invalid certificate matched to a valid one on
every surface property, with per-item certification that the displayed lines
leave the query underdetermined. Adapting the manipulation also changes what a
null means: a model that ignores an injected mistake in its own trace is
unfaithful, whereas a model that follows a corrupted external certificate is
unsafe.

\subsection{Solver-augmented and tool-augmented reasoning}

Neurosymbolic pipelines put the formal step outside the model. Logic-LM
translates a problem into a task-specific symbolic language, calls a solver, and
refines the formalisation from solver error messages \citep{pan2023logiclm};
LINC has the model act as a semantic parser into first-order logic and offloads
deduction to a theorem prover \citep{olausson2023linc}; SatLM generates a
declarative specification and dispatches it to an automated theorem prover
\citep{ye2023satlm}; PAL emits a program and defers execution to an interpreter
\citep{gao2022pal}. Related work extends the pattern to self-supervised API use
\citep{schick2023toolformer}, to retrieval-augmented premise selection in
interactive theorem proving \citep{yang2023leandojo}, and to retrieval as a
general source of external evidence \citep{lewis2020retrievalaugmented}.

These systems report end-to-end accuracy and attribute the gain to the pipeline.
Two things are consequently left open. First, the contribution of the
\emph{reasoning state} is not separated from the contribution of the
\emph{answer}: because the solver's output normally contains both, no reported
number distinguishes a model that reasons over supplied premises from one that
reads off a supplied conclusion. Second, the solver is assumed correct, so the
behaviour of the pipeline under a wrong-but-well-formed solver output is not
measured. Adjacent work shows the model side is not neutral to what the context
contains: irrelevant material inserted into a problem degrades accuracy
\citep{shi2023large}, and models are highly receptive to coherent external
evidence that contradicts their parametric knowledge \citep{xie2023adaptive}.
This paper decomposes the solver's contribution into an answer channel and a
state channel, measures each on matched arms, and adds an arm in which the
certificate is valid-looking but points the wrong way. The state channel carries
real validity information; it supplies no protection against a corrupted
apparatus.

\subsection{Synthetic deductive reasoning benchmarks}

The task family used here descends from rule-based deductive datasets stated in
natural language. RuleTaker showed that transformers can assign truth values to
candidate implications of a theory expressed in language and generalise to
deeper chaining than seen in training \citep{clark2020transformers}. ProofWriter
added generated proofs and open-world semantics with an explicit unknown class
\citep{tafjord2020proofwriter}. PrOntoQA generates items from a synthetic world
model in first-order logic so that a generated chain can be parsed into a
symbolic proof and analysed step by step, and finds that individual deduction
steps are usually correct while proof planning is not
\citep{saparov2022language}. FOLIO supplies human-written premises with verified
first-order logic annotations \citep{han2022folio}. Two properties of these
resources matter for the present design: order effects are known to be large in
deductive prompting \citep{chen2024premise}, and the sets are public and predate
the models evaluated on them, so contamination cannot be excluded --- a problem
documented for modern benchmarks generally
\citep{deng2024investigating,golchin2023time} and sharpened by the finding that
performance can drop sharply on counterfactual variants of familiar tasks
\citep{wu2023reasoning}.

The gap is that these benchmarks supply items and proofs but not the matched
adversarial certificate variants an attribution question needs. This paper
generates items after model release with per-item nonce vocabularies, and
constructs the matched invalid, misleading, order-scrambled and depth-truncated
certificates that the public sets do not provide. A public open-world corpus is
still used, but only as an independent check on the forward-closure certifier
that labels the underdetermined items; no item from it is scored by any model
here.

\subsection{Response bias and the limits of accuracy}

Prompted classification is well documented to be biased rather than merely
noisy. \citet{zhao2021calibrate} show that few-shot accuracy swings from near
chance to near state of the art with prompt format and example order because the
model carries a prior toward particular answer strings, and propose estimating
and removing that prior. Option order alone moves multiple-choice accuracy
substantially \citep{pezeshkpour2023large}, and the effect traces to a token-level
prior over option identifiers rather than to the content of the options
\citep{zheng2023large}. Scoring by raw likelihood introduces a further
confound, since strings compete for probability mass independently of their
correctness \citep{holtzman2021surface}. On the abstention side, models are
partially able to assess whether their own answers are right
\citep{kadavath2022language}, and can be evaluated on whether they recognise
questions that admit no answer \citep{yin2023large}.

What this literature diagnoses, evaluation practice in solver-augmented
reasoning has not absorbed: bare accuracy remains the reported quantity. On a
balanced panel, a degenerate single-label responder is indistinguishable from a
chance-level competent one, and both are indistinguishable from a genuine
partial discriminator with a shifted criterion. Signal detection theory
separates exactly these
\citep{macmillan1991detection,stanislaw1999calculation}, and this paper
therefore reports $d'$ and criterion $c$ alongside accuracy for every arm, with
paired significance from the exact McNemar test \citep{mcnemar1947note} and
interval estimates from the bias-corrected accelerated bootstrap
\citep{efron1987better}. Two concrete cases motivate the practice rather than
illustrate it: a baseline arm scoring $50.0\%$ on a balanced panel turns out to
have $d' = 0.00$, and a prespecified qualification criterion based on a single
class recall was passed by models that answer one label almost everywhere. Both
were caught only because sensitivity was reported separately from accuracy; one
of the two is a criterion of this paper's own, replaced post hoc and labelled as
such.

\section{Method}

\subsection{The answer-evidence decomposition}
\label{sec:decomposition}

A solver certificate carries two things at once: the \emph{answer} to the query,
and the \emph{state} from which that answer follows. End-to-end accuracy
conflates them. We separate them by arranging arms as a $2\times2$ of answer
evidence by proof state, with a no-material baseline outside the square.

Let \textsc{irrelevant} supply a valid derivation for a different entity (no
answer, no relevant state), \textsc{conclusion\_only} the derived terminal
literal alone (answer, no state), \textsc{proof\_prefix} the chain with that
literal withheld (state, no answer), and \textsc{full} both. A contrast along an
\emph{edge} of this square moves one factor; a contrast along the
\emph{diagonal} moves two, and a null on it attributes to neither. This is not a
statistical subtlety but an identification one: the prespecified primary
contrast of the experiment we reanalyse, \textsc{proof\_prefix} minus
\textsc{conclusion\_only}, is exactly that diagonal.

Reading the four edges instead gives a state main effect, the mean of the two
edges along which state is added, and an answer main effect defined
symmetrically.

\subsection{Surface-matched invalidity}
\label{sec:broken}

A prefix that stops one step short is not a clean state manipulation. It still
names the query's subject and predicate in adjacent lines, so surface
co-occurrence remains available. To isolate validity we construct
\textsc{broken\_chain}: one intermediate link is cut so the displayed lines no
longer reach the query, while every surface property is held fixed.

Concretely, against the matched \textsc{truncate\_1} arm we preserve the
query-entity occurrence count (audited per item, not approximately), the
presence of the query predicate in the final rule, the line count, the exact
fact/rule/literal type sequence, and token length within a recorded tolerance.
Every \emph{rule} displayed is a real rule of the theory; exactly one derived
literal is false. All 192 items are certified by forward-closure saturation
(\cref{sec:closure}) to leave the query \emph{undetermined} under the displayed
lines. The decomposition follows:
\begin{align}
\text{total state effect} &= \textsc{truncate\_1} - \textsc{irrelevant},\\
\text{surface component} &= \textsc{broken\_chain} - \textsc{irrelevant},\\
\text{validity component} &= \textsc{truncate\_1} - \textsc{broken\_chain}.
\end{align}

Four companion arms separate the remaining confounds.
\textsc{same\_entity\_irrelevant} supplies a valid derivation about the
\emph{query} entity toward an unrelated predicate: it names the query subject
four times where \textsc{irrelevant} names it zero times, and neither contains
the query predicate, so the pair isolates entity repetition alone.
\textsc{shuffled} presents \textsc{truncate\_1}'s lines in seeded random order,
separating a chain from a bag of statements. \textsc{truncate\_2} and
\textsc{truncate\_3} withhold two and three steps, giving a depth ladder.
\textsc{misleading} supplies a valid-looking chain whose single fabricated final
rule establishes the query's negation; the ground-truth answer is unchanged, so
a model that follows the supplied state answers incorrectly.

\subsection{Certifying underdetermination}
\label{sec:closure}

Showing that a query \emph{is} derivable is easy: exhibit the proof. Showing it
is \emph{not} requires saturating the theory's forward closure. Our certifier
operates under open-world semantics with explicit negation, so a negative
literal is derivable only if something derives it, never by failure to prove the
positive. Rule antecedents are ordinary literals, with no negation as failure,
which makes naive forward chaining sound and complete on this fragment.

Two safeguards matter for generated theories. The chaining loop carries an
explicit round cap and reports a \texttt{saturated} flag, so a malformed theory
fails loudly rather than spinning; and any theory deriving both an atom and its
negation is marked inconsistent and its items rejected rather than labelled. An
item is \emph{undetermined} when neither the query nor its negation appears in
the closure.

We validated the certifier before generating any panel, against the open-world
splits of an independently labelled corpus \citep{tafjord2020proofwriter} at
depths $0$, $1$, $2$, $3$ and $5$: $23{,}240$ of $23{,}240$ questions agree,
$100.00\%$, with perfectly diagonal confusion and zero unusable theories,
including $10{,}440$ undetermined items. No item enters a sealed panel without a
certification record.

\subsection{Measurement: sensitivity, not accuracy}
\label{sec:measurement}

Accuracy on a balanced panel is uninterpretable when the responder is biased.
One arm below scores $50.0\%$ --- apparent chance --- with an identical
$9/96$ ``Yes'' rate in \emph{both} classes, that is $d' = 0.00$: no
discriminative signal whatever, at or below the trivial always-one-label
strategy. We therefore report sensitivity $d'$ and criterion $c$ alongside
accuracy for every arm \citep{macmillan1991detection,stanislaw1999calculation},
using the loglinear correction for extreme rates. For the three-class panel the
headline scalar is balanced accuracy, the mean of per-class recalls, under which
chance is $33.3\%$ and a single-label responder scores exactly chance.

Paired contrasts report the complete $2\times2$ table, the exact two-sided
McNemar test \citep{mcnemar1947note} as the decision criterion, and a seeded
$100{,}000$-resample paired bias-corrected accelerated bootstrap interval
\citep{efron1987better}, with Holm adjustment within declared families.

\section{Experimental Setup}

\subsection{Task, panels and scoring}

All three experiments use rule-chaining queries over small theories of facts and
universally quantified rules. \Cref{tab:panels} lists the panels. The three
sealed panels are procedurally generated after model release with per-item nonce
vocabularies, are mutually disjoint, and fix derivation depth at four by
construction, recorded per item. The public corpus is used \emph{only} to
validate the certifier; no item from it is scored by any model here.

\begin{table*}[t]
\centering\small
\caption{Datasets. The three sealed panels are nonce-generated and disjoint; the
public corpus validates the certifier only.}
\label{tab:panels}
\begin{tabular}{llrl}
\toprule
Panel & Classes & Items & Used by \\
\midrule
\texttt{b14} & entailed / contradicted & 192 & Experiment 1 \\
\texttt{b15} & entailed / contradicted & 192 & Experiment 2 \\
\texttt{b16} & entailed / contradicted / undetermined & 192 & Experiment 3 \\
ProofWriter OWA \citep{tafjord2020proofwriter} & three-way & 23{,}240 & certifier validation \\
\bottomrule
\end{tabular}
\end{table*}

Responses are scored by direct next-token likelihood over verified single-token
candidates: \texttt{Yes}/\texttt{No} for Experiments 1 and 2,
\texttt{Yes}/\texttt{No}/\texttt{Unknown} for Experiment 3. The model generates
no free text, so nothing is parsed and no verdict is inferred from prose. Every
prompt is four lines --- instruction, theory, solver record, query --- of which
only the solver record varies across arms; the theory is visible in all arms.

\subsection{Models}

Experiments 1 and 2 use the same frozen 3B instruction-tuned checkpoint at
4-bit precision \citep{team2024qwen}, pinned to byte-identical weights, so the
anchors of the first transfer to the second without a weights caveat.
Experiment 3 scores ten instruction-tuned models from $0.5$B to $7.6$B
\citep{team2024qwen,grattafiori2024llama,team2025gemma,microsoft2025phi,team2024olmo,allal2025smollm},
all at bfloat16 so precision is not confounded with model identity. An
eligibility screen confirmed that all ten load, complete a forward pass, and
tokenise all three candidates as single distinct tokens.

\subsection{Sealed protocol}

Panels are hash-stamped before scoring and the runner refuses a modified panel.
Each response is written to a receipt carrying a SHA-256 of its own contents,
and the answer authority is opened only once every receipt exists, so nothing
can be tuned to a result. Audits depend on the standard library alone, so
verification never requires the environment to resolve. In total $12{,}480$
prompts were scored with zero receipt digest failures.

Experiment 3 makes a deliberate protocol change: \emph{no gate blocks the run}.
Where earlier practice excluded a model that failed a delivery screen, delivery
is itself the capability under test here, so a model that cannot emit a verdict
reliably is a finding about substrate viability rather than a reason to drop it.
Tooling failures --- a model that will not load, or whose tokeniser cannot
represent the candidates --- remain exclusions and are not findings; the screen
found none.

\section{Results}

\subsection{Reading the factorial as a \texorpdfstring{$2\times2$}{2x2}}

\Cref{fig:factorial} shows the five arms of Experiment 1 arranged by
\cref{sec:decomposition}. The prespecified primary contrast,
$\textsc{proof\_prefix} - \textsc{conclusion\_only}$, was $-4.7$pp with exact
McNemar $p = 0.078$ and was reported as a null. It is the diagonal.

\begin{figure*}[t]
\centering
\includegraphics[width=\columnwidth]{figures/answer_state_factorial.png}
\caption{The five arms form a $2\times2$ of answer evidence by proof state, with
the no-material baseline outside it. Each cell gives correct items of 192.
Arrows mark the four one-factor edges; the dashed diagonal is the prespecified
primary, which moves both factors at once and therefore attributes its null to
neither.}
\label{fig:factorial}
\end{figure*}

The four edges are reported in \cref{tab:edges}. \textbf{This reading is post
hoc}: it is a different analysis of sealed data, not a re-run, and every value
is re-derived from the 960 raw receipts by a script that reconstructs the answer
key and asserts it reproduces all five published per-arm counts before
reporting. The state main effect is $+22.4$pp and the answer main effect
$+27.1$pp; the answer literal retains only $30\%$ of its value once state is
present.

\begin{table*}[t]
\centering\small
\caption{The four one-factor edges of the $2\times2$, and the prespecified
diagonal. Post hoc reanalysis of sealed data.}
\label{tab:edges}
\begin{tabular}{lrrl}
\toprule
Contrast & Effect & Exact McNemar $p$ & Reported originally \\
\midrule
Add answer, no state present   & $+41.67$pp & $1.7\times10^{-24}$ & no \\
Add state, no answer present   & $+36.98$pp & $1.5\times10^{-19}$ & secondary \\
Add answer, state present      & $+12.50$pp & $1.2\times10^{-7}$  & no \\
Add state, answer present      & $+7.81$pp  & $6.1\times10^{-5}$  & secondary \\
\midrule
\emph{Diagonal (prespecified primary)} & $-4.69$pp & $0.078$ & primary \\
\bottomrule
\end{tabular}
\end{table*}

Splitting by class overturns the baseline description. The model answers
\texttt{No} to $88\%$ of items. \Cref{tab:sdt} gives the signal detection
summary: \textsc{irrelevant} returns an identical $9/96$ ``Yes'' rate in both
classes, $d' = 0.00$. A trivial always-\texttt{No} strategy scores $96/192$;
\textsc{none} scored 91 and \textsc{irrelevant} 96, at or below it. The
contradicted class sits at ceiling in every certificate arm and contributed
$0$ of the $21$ discordant pairs in the primary contrast, so the effective
sample was 96, not 192.

\begin{table*}[t]
\centering\small
\caption{Signal detection for Experiment 1. Apparent chance and zero
sensitivity are not the same thing.}
\label{tab:sdt}
\begin{tabular}{lrrrr}
\toprule
Arm & Hits /96 & False alarms /96 & $d'$ & Criterion $c$ \\
\midrule
\textsc{none}            &  9 & 14 & $-0.26$ & $+1.17$ \\
\textsc{irrelevant}      &  9 &  9 & $\mathbf{0.00}$ & $+1.29$ \\
\textsc{conclusion\_only}& 80 &  0 & $+3.52$ & $+0.81$ \\
\textsc{proof\_prefix}   & 71 &  0 & $+3.20$ & $+0.97$ \\
\textsc{full}            & 95 &  0 & $+4.72$ & $+0.20$ \\
\bottomrule
\end{tabular}
\end{table*}

\subsection{Validity versus surface overlap}

Experiment 2 holds surface form fixed and destroys validity alone. The primary
contrast, evaluated on the entailed subset ($n = 96$) as declared before the run
because the contradicted class contributes no discordance, is
$\textsc{truncate\_1} - \textsc{broken\_chain} = +60.4$pp, exact McNemar
$p = 6.9\times10^{-18}$, BCa $[+49.0, +68.8]$, with a paired table of $1$ both
correct, $58$ \textsc{truncate\_1}-only, $0$ \textsc{broken\_chain}-only and
$37$ neither. Sensitivity moves with it: $d'$ drops by $+2.446$, satisfying the
same-sign requirement declared in advance, so this is a change in
discrimination and not in criterion placement.

The decomposition is therefore a total state effect of $+61.5$pp comprising a
surface component of $+1.0$pp ($p = 1$) and a validity component of $+60.4$pp:
a \textbf{validity share of $98.3\%$}. \Cref{fig:arms} and \cref{tab:arms} give
the full arm set.

\begin{figure*}[t]
\centering
\includegraphics[width=0.86\textwidth]{figures/arm_decomposition.png}
\caption{Ten matched arms on one item set, differing only in the solver record.
Left: correct on the entailed subset (of 96). Right: sensitivity $d'$, signed.
Destroying validity alone removes almost the whole effect. The
\textsc{misleading} arm, whose fabricated final rule points the wrong way, is
followed to the wrong answer on 189 of 192 items.}
\label{fig:arms}
\end{figure*}

\begin{table*}[t]
\centering\small
\caption{Experiment 2, all ten arms. Five arms return identical results,
including \textsc{same\_entity\_irrelevant}, which names the query entity four
times.}
\label{tab:arms}
\begin{tabular}{lrrr}
\toprule
Arm & All /192 & Entailed /96 & $d'$ \\
\midrule
\textsc{none}                    &  96 &  0 & $0.000$ \\
\textsc{irrelevant}              &  96 &  0 & $0.000$ \\
\textsc{same\_entity\_irrelevant}&  96 &  0 & $0.000$ \\
\textsc{truncate\_3}             &  96 &  0 & $0.000$ \\
\textsc{truncate\_2}             &  96 &  0 & $0.000$ \\
\textsc{truncate\_1}             & 155 & 59 & $2.853$ \\
\textsc{broken\_chain}           &  97 &  1 & $0.407$ \\
\textsc{shuffled}                & 110 & 14 & $1.527$ \\
\textsc{misleading}              &   3 &  0 & $-4.363$ \\
\textsc{full}                    & 192 & 96 & $5.131$ \\
\bottomrule
\end{tabular}
\end{table*}

Three supporting contrasts sharpen the reading. Entity repetition explains
nothing: $\textsc{same\_entity\_irrelevant} - \textsc{irrelevant} = +0.0$pp
($p = 1$), despite the former naming the query subject four times and the latter
not at all. Order matters: $\textsc{truncate\_1} - \textsc{shuffled} = +46.9$pp
($p = 6.8\times10^{-13}$), so the model is not reading a bag of statements.
And depth is a cliff, not a slope:
$\textsc{truncate\_2} - \textsc{truncate\_1} = -61.5$pp
($p = 3.5\times10^{-18}$), with \textsc{truncate\_2} scoring identically to
supplying no certificate at all. Five arms are indistinguishable ---
$96/192$, $0/96$ entailed, $d' = 0.000$, $c = 2.565$ --- so behaviour is binary:
either the supplied chain reaches one step from the answer, or the model is
blind to it.

\subsection{Following a corrupted apparatus}

The \textsc{misleading} arm supplies a chain identical to \textsc{full} except
that its final rule, which is fabricated and absent from the theory, flips
polarity. The ground-truth answer is unchanged. The model scores $3$ of $192$,
with $d' = -4.363$: $\textsc{misleading} - \textsc{full} = -100.0$pp
($p = 2.5\times10^{-29}$).

The model therefore follows a valid-looking certificate that points the wrong
way on 189 of 192 items. Read against the $98.3\%$ validity share, the two
results are not in tension but jointly informative: the model demonstrably
tracks whether a supplied derivation connects, and that tracking provides no
defence whatever when the derivation connects to the wrong conclusion. Validity
tracking is not verification.

\subsection{Which models can serve as the interface}

Experiment 3 scores ten models on the three-class panel. We declared in advance
that a model relays non-determination if it reaches $\geq 80\%$ recall on
undetermined items with full solver material. \textbf{Eight of ten passed, and
the criterion was invalid.} It measures recall on a single class, which a model
answering \texttt{Unknown} to everything maximises while discriminating nothing
--- and four models do exactly that, emitting \texttt{Unknown} on $62$--$96\%$
of all items, three of them never emitting \texttt{No} at all. This is the
same failure the $d' = 0.00$ arm exhibits in Experiment 1, mirrored onto a
different label.

We replace it \textbf{post hoc} with a floor on every class: a model is a viable
substrate if its minimum per-class recall is at least $0.50$. Collapse cannot
game this. Both verdicts are retained in the released results;
the superseded one is not deleted. \Cref{tab:models} and \cref{fig:size} report
the outcome.

\begin{figure}[t]
\centering
\includegraphics[width=\columnwidth]{figures/substrate_size_curve.png}
\caption{Balanced accuracy for ten models with no solver material (open marker)
and with full material (filled), ordered by parameter count. Chance is $33.3\%$.
Colour gives the corrected viability verdict. Three of ten are viable; below
1.7B the interface collapses to a single label whatever the solver supplies.}
\label{fig:size}
\end{figure}

\begin{table*}[t]
\centering\small
\caption{Experiment 3. Balanced accuracy (\%), chance $=33.3$. The prespecified
criterion passed eight of ten; the corrected criterion passes three.}
\label{tab:models}
\begin{tabular}{lrrrrrl}
\toprule
Model & Params (B) & \textsc{none} & \textsc{full} & Min recall & Max label share & Verdict \\
\midrule
Qwen2.5-0.5B  & 0.5 & 35.4 & 63.5 & 0.0  & 61.5 & collapsed \\
OLMo-2-1B     & 1.0 & 35.9 & 42.2 & 0.0  & 91.1 & collapsed \\
Llama-3.2-1B  & 1.2 & 33.9 & 37.0 & 0.0  & 95.3 & collapsed \\
Qwen2.5-1.5B  & 1.5 & 33.3 & 37.0 & 0.0  & 96.4 & collapsed \\
SmolLM2-1.7B  & 1.7 & 33.3 & 69.3 & 37.5 & 44.8 & below floor \\
Qwen2.5-3B    & 3.0 & 41.1 & \textbf{90.6} & 71.9 & 42.7 & \textbf{viable} \\
Llama-3.2-3B  & 3.2 & 45.8 & 80.7 & 42.2 & 49.5 & below floor \\
Phi-4-mini    & 3.8 & 44.8 & \textbf{90.6} & 71.9 & 42.7 & \textbf{viable} \\
Gemma-3-4B    & 4.3 & 63.5 & \textbf{96.4} & 89.1 & 36.5 & \textbf{viable} \\
Qwen2.5-7B    & 7.6 & 33.3 & 71.4 & 14.1 & 62.0 & collapsed \\
\bottomrule
\end{tabular}
\end{table*}

Three findings follow. Supplying full solver material takes viable models from
at or near chance to $90$--$96\%$ balanced accuracy on a task that includes
abstention, so the architecture works above a floor. That floor is near 3B:
below 1.7B, models collapse to a single label whatever the solver supplies, and
an apparatus cannot help an interface that cannot express three outcomes. And
the effect is not monotonic in scale --- the 7.6B model collapses ($62.0\%$
\texttt{Unknown}, minimum per-class recall $14.1$) while a 3B model does not.
Model is not a randomised factor here, so these comparisons are descriptive and
we make no significance claim about scale.

\section{Limitations}

\subsection{Scope of the mechanism claim}

The inference we measure is \emph{exactly one step deep}. \textsc{truncate\_2}
scores identically to supplying no certificate, so ``validity tracking'' should
be read as: the model can complete a final inference when both premises are
present and adjacent, and can detect when the chain that would license it is
broken. It cannot chain two steps. That is a hard bound on how this result may
be described.

The full theory is visible in every arm, so the query remains derivable from the
theory regardless of what the certificate says. Breaking a certificate is
diagnostic only because the model scores $0/96$ working from the theory alone. A
model that used a broken certificate as a pointer back into the theory would be
performing inference of a different kind; our design does not exclude that
reading, and we record it as an interpretation caveat rather than a controlled
alternative.

\textsc{truncate\_2} and \textsc{truncate\_3} necessarily drop the query
predicate out of the certificate, since it lives in the final rule. The depth
ladder therefore confounds depth with predicate presence and must be read
against \textsc{same\_entity\_irrelevant}, which holds entity repetition high
while removing the predicate.

\subsection{Post hoc analyses and superseded criteria}

Two analyses in this paper are post hoc and are labelled where they are used.
The $2\times2$ reading of Experiment 1 is a different analysis of sealed data,
not a re-run. The viability criterion of Experiment 3 replaced a prespecified
threshold that the data showed to be invalid; the original is retained in the
released results.

Experiment 1 also carries design errors we report rather than repair. It
prespecified a $15$pp primary effect while its comparison arm reached $91.7\%$,
leaving $8.3$pp of headroom, so the target was arithmetically unreachable. It
recorded no proof depth and so cannot stratify its failures. And its task
identifiers encoded the semantic class, so its answer authority was
reconstructible from the public panel --- convenient for our reanalysis, but not
a sealed authority. Experiments 2 and 3 fix all three.

\subsection{Generalisation}

This is one task family, small models only, and synthetic panels; the mechanism
result rests on a single model, and the ten-model sweep addresses substrate
viability rather than mechanism. Nonce vocabularies establish item novelty, not
independence from all relevant pretraining patterns
\citep{golchin2023time,deng2024investigating}.

An undetermined \textsc{proof\_prefix} cannot contain the query predicate,
because the solver has nothing to say about it, whereas entailed and
contradicted prefixes can. Undetermined items therefore offer fewer surface
cues, and a surface-matching model will look worse on them for reasons unrelated
to abstention; per-arm entity-mention counts are recorded so this is analysable
rather than discovered later.

Finally, every theory here arrives \emph{already formalised}. Nothing in this
paper tests whether a model can take a framework described in prose and produce
the solver's input, which is the step on which any claim of working in an
unfamiliar framework depends. \Cref{sec:closure} and the corruption result make
that gap consequential rather than merely open: a formalisation error would not
be caught downstream, but followed.

\section{Conclusion}

Solver assistance improves a small model's accuracy on rule-chaining queries,
and the improvement is not text matching. Holding line count, line types,
entity occurrence count, predicate presence and token length fixed and
destroying validity alone removes $98.3\%$ of the state effect. Entity
repetition, the most plausible surface confound, contributes nothing measurable.

Two qualifications travel with that result and should not be separated from it.
The inference is one step deep, so this is a narrow competence and not a general
reasoning claim. And it provides no protection against a wrong apparatus: the
same model that detects a broken chain follows a fabricated one to the wrong
answer on 189 of 192 items. For architectures that externalise reasoning to a
verified component and use the model as the interface, that asymmetry is the
operationally important finding --- validity tracking is not verification, and a
system built on it inherits the apparatus's errors in full.

The measurement lessons generalise beyond this task. Twice in this programme a
number that looked like competence was a degenerate responder: an arm at
$50.0\%$ accuracy with $d' = 0.00$, and a prespecified criterion that a model
answering \texttt{Unknown} to everything maximised. Both were caught only by
reporting sensitivity alongside accuracy. We recommend the practice for any
evaluation of tool-augmented systems, together with the answer-evidence
decomposition, which requires no new data collection --- only that the arms be
arranged so that each contrast moves one factor at a time.

\section*{Reproducibility Statement}

\textbf{Code and data.} The supplementary material accompanying this submission
is self-contained: sealed panels and answer authorities, all $12{,}480$
per-response receipts, the generators, runners, audits and analysis scripts, and
the forward-closure certifier with its validation harness. Every number in this
paper can be re-derived from a fresh unpack with no network access, no model
weights and no external packages --- the audits use the Python standard library
only, so verification does not depend on resolving an environment. A permanent
archive with a DOI will accompany the camera-ready version.

All panels, answer authorities, per-response receipts, audit scripts and
analysis scripts are released. Every number in this paper is re-derivable from
the receipts without loading a model: the audits use the standard library only,
so verification does not depend on resolving an environment. Panels are
hash-stamped and the runners refuse a modified panel; each receipt carries a
SHA-256 of its own contents, and all $12{,}480$ verified with zero failures.
Model weights are not redistributed, but each run records the pinned revision
and a hash of the weight files before those files are deleted. The
forward-closure certifier and its validation against the public corpus are
included.

\section*{AI-Use Statement}

A large language model was used as a coding and drafting assistant throughout:
generating experiment code, analysis scripts and figure code, and drafting this
manuscript. All experimental design decisions, the interpretation of results,
and the claim boundaries were made and verified by the authors. Every numeric
value reported was produced by the released analysis scripts from sealed
receipts, not written by a model; the released audits allow any reader to
re-derive them independently.

\bibliographystyle{plainnat}
\bibliography{refs}

\end{document}
